\documentclass[11pt]{article}

\usepackage[margin=1in]{geometry}
\usepackage{cmap}
\usepackage[T1]{fontenc}
\usepackage[utf8]{inputenc}
\usepackage{booktabs}
\usepackage{longtable}
\usepackage{array}
\usepackage{xcolor}
\usepackage{hyperref}
\usepackage{listings}

\hypersetup{
  hidelinks,
  pdftitle={MATRIX Remote QM Execution on oracle},
  pdfauthor={MATRIX Development Team}
}

\definecolor{MatrixBlue}{HTML}{245C73}
\definecolor{MatrixSlate}{HTML}{4A5561}
\definecolor{MatrixGold}{HTML}{C88A2D}
\definecolor{MatrixBg}{HTML}{F7F8F6}
\emergencystretch=2em

\lstdefinestyle{matrix}{
  basicstyle=\ttfamily\small,
  breaklines=true,
  columns=fullflexible,
  keepspaces=true,
  frame=single,
  rulecolor=\color{MatrixSlate!45},
  backgroundcolor=\color{MatrixBg},
  xleftmargin=0.5em,
  xrightmargin=0.5em
}

\newcommand{\framework}{\texttt{MATRIX}}
\newcommand{\host}{\texttt{oracle}}
\newcommand{\xyzin}{\texttt{xyzin}}

\title{\textbf{MATRIX Remote QM Execution on oracle}\\
\large Running Gaussian, Molpro and ORCA jobs for MATRIX workflows}
\author{MATRIX Development Notes}
\date{\today}

\begin{document}
\maketitle

\begin{abstract}
This manual describes the operational contract for running external quantum
chemistry programs on the private Linux host \host{} and bringing the completed
results back into \framework{}.  The remote machine is used only for execution
of heavy QM jobs.  \framework{} remains the owner of the molecular state,
shared \xyzin{} sections, generated Gaussian GIC inputs and normalized
post-processing through the single QM-adapter layer.
\end{abstract}

\tableofcontents
\newpage

\section{Scope}

The \host{} machine is intended for long background calculations with Gaussian,
Molpro and ORCA.  It is not a replacement for the local \framework{} source
tree and it must not become a second set of parsers.  The normal workflow is:

\begin{enumerate}
  \item prepare or validate the molecule locally with \framework{};
  \item generate the QM input locally when \framework{} owns the input format,
        for example Gaussian ReadAllGIC files from frozen NEO/GICForge data;
  \item copy the input file to \host{};
  \item submit the job in the remote \texttt{\textasciitilde/matrix} work area;
  \item copy completed logs, checkpoint-derived files or Molden files back;
  \item promote the output through the corresponding \framework{} adapter into
        shared \xyzin{} sections.
\end{enumerate}

Passwords, personal credentials and site-specific secrets must not be written
in the repository, in job inputs or in shell scripts.  Use the interactive SSH
prompt or a normal SSH key managed outside the repository.

\section{Remote Directory Layout}

The remote execution area is \texttt{\textasciitilde/matrix}.  The following
subdirectories have fixed meanings:

\begin{longtable}{
  >{\raggedright\arraybackslash}p{0.25\linewidth}
  >{\raggedright\arraybackslash}p{0.65\linewidth}}
\toprule
Directory & Purpose \\
\midrule
\texttt{\textasciitilde/matrix/bin} & MATRIX helper scripts on the remote host. \\
\texttt{\textasciitilde/matrix/inputs} & Input files copied from the local workstation. \\
\texttt{\textasciitilde/matrix/jobs} & Per-job working directories created by the submit helper. \\
\texttt{\textasciitilde/matrix/logs} & Lightweight submit logs and redirected program output. \\
\texttt{\textasciitilde/matrix/outputs} & Optional collection area for selected final products. \\
\texttt{\textasciitilde/matrix/repo} & Optional clone or mirror of MATRIX when needed for tests. \\
\texttt{\textasciitilde/matrix/scratch} & MATRIX-owned scratch area; large Gaussian scratch is controlled separately. \\
\bottomrule
\end{longtable}

Gaussian scratch is configured through \texttt{GAUSS\_SCRDIR}.  The current
policy is to keep large Gaussian scratch under the local scratch filesystem
configured on \host{}, not inside the Git repository.

\section{Environment}

The login shell on \host{} sources \texttt{\textasciitilde/.matrix\_env.sh}
from both \texttt{\textasciitilde/.bashrc} and, when available,
\texttt{\textasciitilde/.zshrc}.  The environment file is intentionally narrow:
it exposes only the production Gaussian choices and the supported module stack.

\begin{lstlisting}[style=matrix]
setgau gdv32
setgau g16
module load molpro/2025.3
module load orca/6.1.1
\end{lstlisting}

Older GDV variants are disabled for MATRIX use.  Jobs must use only
\texttt{gdv32} or \texttt{g16}.  Do not source historical external-code
environments such as \texttt{elecext-env}, \texttt{External} or
\texttt{RunGauExternal}; they are outside the MATRIX execution contract.

Useful checks after login are:

\begin{lstlisting}[style=matrix]
command -v gdv
command -v g16
command -v molpro
command -v orca
module list
echo "$GAUSS_SCRDIR"
\end{lstlisting}

The expected Molpro version for this machine is Molpro 2025.3.  If
\texttt{module load molpro/2025.3} stops working, treat that as an environment
configuration problem and do not silently fall back to an older Molpro module.

\section{Direct Program Commands}

The direct commands are useful for small checks and for debugging the
environment.

\subsection{Gaussian GDV32}

\begin{lstlisting}[style=matrix]
setgau gdv32
gdv32run molecule.gjf
\end{lstlisting}

Use GDV32 for Gaussian jobs that need the current group-specific Gaussian
features.  For MATRIX-generated ReadAllGIC optimizations, the Gaussian route
must use \texttt{Opt=ReadAllGIC}; for single-point jobs with a GIC geometry
read, use the corresponding \texttt{Geom} form.  MATRIX symmetrizes the GICs
itself, so Gaussian-side automatic GIC symmetrization should not be requested.

\subsection{Gaussian G16}

\begin{lstlisting}[style=matrix]
setgau g16
g16run molecule.gjf
\end{lstlisting}

Use G16 when the job does not need GDV32-specific behavior.  The same MATRIX
adapter policy applies: logs, FCHK files and QFF data are copied back and then
promoted through \texttt{matrix gaussian ...} commands.

\subsection{Molpro}

\begin{lstlisting}[style=matrix]
module load molpro/2025.3
molpro molecule.com
\end{lstlisting}

Molpro outputs enter MATRIX through the Molpro adapter.  Downstream tools
should consume the promoted \xyzin{} sections and must not parse Molpro text
privately.

\subsection{ORCA}

\begin{lstlisting}[style=matrix]
module load orca/6.1.1
orca molecule.inp > molecule.out
\end{lstlisting}

ORCA calculations can be launched on \host{} for electronic-structure data and
future MATRIX adapters.  Until a given ORCA output quantity is normalized into
a shared \xyzin{} section, scientific MATRIX tools should not depend on
tool-local ORCA parsing.

\section{Background Submission}

For production use, prefer the lightweight MATRIX submit helpers installed in
\texttt{\textasciitilde/matrix/bin}.  The source version of the submit helper
is stored in the MATRIX repository as
\path{scripts/remote/oracle/matrix-submit}.  It creates one working
directory per job, records the native program output, exposes a stable log
link in \texttt{\textasciitilde/matrix/logs} and keeps a small metadata file
with the command and process id.

\begin{lstlisting}[style=matrix]
matrix-submit gdv32 ~/matrix/inputs/molecule.gjf
matrix-submit g16 ~/matrix/inputs/molecule.gjf
matrix-submit molpro ~/matrix/inputs/molecule.com
matrix-submit orca ~/matrix/inputs/molecule.inp
\end{lstlisting}

Monitor jobs with:

\begin{lstlisting}[style=matrix]
matrix-status
matrix-tail 20260628-101010-gdv32-molecule
matrix-tail ~/matrix/logs/20260628-101010-gdv32-molecule.log
\end{lstlisting}

\texttt{matrix-submit} is a background wrapper based on ordinary Unix process
management.  It is not a queueing system.  It is appropriate for this private
machine, but it does not implement resource accounting, priorities or node
scheduling.  Each job writes \texttt{metadata.txt}; the fields
\texttt{log}, \texttt{native\_output} and \texttt{stdout} distinguish the
stable log link, the scientific program output in the work directory and any
launcher stdout/stderr stream.

\section{Local MATRIX Remote Commands}

The preferred local interface is the \framework{} remote QM wrapper.  It uses
SSH and SCP only as transport, calls the same remote
\texttt{matrix-submit} helper and never stores credentials.

\begin{lstlisting}[style=matrix]
matrix qm remote-submit molecule.gjf --engine gdv32 --host enzo@oracle
matrix qm remote-submit molecule.com --engine molpro --host enzo@oracle
matrix qm remote-submit molecule.inp --engine orca --host enzo@oracle

matrix qm remote-status --host enzo@oracle
matrix qm remote-fetch JOB_NAME --host enzo@oracle --dest runs
\end{lstlisting}

\texttt{remote-fetch} reads the remote \texttt{metadata.txt}, copies the
\texttt{native\_output} file into a local job directory and writes
\texttt{remote\_qm\_fetch\_manifest.json}.  When a safe adapter is already
defined, fetching and promotion can be combined:

\begin{lstlisting}[style=matrix]
matrix qm remote-fetch JOB_NAME --host enzo@oracle \
  --dest runs --promote molpro --xyzin molecule.xyzin

matrix qm remote-fetch JOB_NAME --host enzo@oracle \
  --dest runs --promote gaussian-log-hessian --xyzin molecule.xyzin
\end{lstlisting}

\texttt{--promote auto} promotes Molpro outputs because the Molpro geometry
adapter is unambiguous.  Gaussian promotion is deliberately explicit because a
single log can contain different useful sections; choose
\texttt{gaussian-log-hessian}, \texttt{gaussian-rovib},
\texttt{gaussian-electronic} or \texttt{gaussian-fchk}.  ORCA outputs are
fetched and recorded, but not promoted until an ORCA adapter for a specific
\xyzin{} section is defined.

\section{Copying Files}

Manual copying remains a fallback.  Copy inputs from the local workstation to
the remote input directory:

\begin{lstlisting}[style=matrix]
scp molecule.gjf enzo@oracle:~/matrix/inputs/
scp molecule.com enzo@oracle:~/matrix/inputs/
scp molecule.inp enzo@oracle:~/matrix/inputs/
\end{lstlisting}

Submit remotely:

\begin{lstlisting}[style=matrix]
ssh enzo@oracle 'matrix-submit gdv32 ~/matrix/inputs/molecule.gjf'
ssh enzo@oracle 'matrix-submit molpro ~/matrix/inputs/molecule.com'
\end{lstlisting}

After the run, copy back the files needed by MATRIX adapters.  The safest
target is the \texttt{native\_output} path recorded in
\texttt{metadata.txt}; for Gaussian this is normally a \texttt{.log}, for
Molpro a \texttt{.out} and for ORCA a \texttt{.out}.

\begin{lstlisting}[style=matrix]
scp enzo@oracle:~/matrix/jobs/JOB_NAME/molecule.log .
scp enzo@oracle:~/matrix/jobs/JOB_NAME/molecule.fchk .
scp enzo@oracle:~/matrix/jobs/JOB_NAME/molecule.out .
\end{lstlisting}

Use the actual job directory name printed by \texttt{matrix-submit} or listed
by \texttt{matrix-status}.  Keep large scratch files remote unless they are
needed for a specific adapter or viewer.

\section{Returning Results To MATRIX}

The completed QM output must be normalized before GF/PED, MORPHEUS,
VPT2/VCI, DVR, spectroscopy or GUI tools consume it.  Typical local commands
after copying files back are:

\begin{lstlisting}[style=matrix]
matrix gaussian promote-fchk molecule.fchk molecule.xyzin
matrix gaussian promote-log-hessian molecule.log molecule.xyzin
matrix gaussian promote-rovib molecule.log molecule.xyzin
matrix gaussian promote-electronic molecule.log molecule.xyzin
matrix molpro promote-output molecule.out molecule.xyzin
\end{lstlisting}

For Gaussian ReadAllGIC optimization checks, keep the coordinate definition and
the completed Gaussian result separate:

\begin{itemize}
  \item the frozen MATRIX \texttt{\#GIC} section is the authoritative
        coordinate source;
  \item the Gaussian input is only an external representation of those GICs;
  \item the Gaussian log supplies the final optimized geometry, Hessian and
        normal modes;
  \item GF/PED must rebuild B rows from the frozen MATRIX GIC definition on
        the final log geometry.
\end{itemize}

This avoids accidentally trusting a hand-edited or Gaussian-modified input
block as the scientific coordinate definition.

\section{Troubleshooting}

\begin{longtable}{
  >{\raggedright\arraybackslash}p{0.33\linewidth}
  >{\raggedright\arraybackslash}p{0.57\linewidth}}
\toprule
Problem & Checks \\
\midrule
Gaussian command not found &
Run \texttt{setgau gdv32}, then \texttt{command -v gdv}.  For standard Gaussian
run \texttt{setgau g16}, then \texttt{command -v g16}. \\
Wrong Gaussian executable &
Use only \texttt{gdv32} or \texttt{g16}.  Do not use older GDV aliases. \\
Molpro unavailable &
Run \texttt{module load molpro/2025.3} and \texttt{molpro --version}. \\
ORCA unavailable &
Run \texttt{module load orca/6.1.1} and \texttt{command -v orca}. \\
No live output &
Use \texttt{matrix-tail JOB\_NAME} or inspect the log in
\texttt{\textasciitilde/matrix/logs}. \\
Possible scratch/disk issue &
Check \texttt{df -h \textasciitilde{} /local/scratch} and the program log. \\
MATRIX tool cannot use a result &
Promote the QM output first and verify that the required \xyzin{} section is
present with \texttt{matrix validate} or the GUI readiness panel. \\
\bottomrule
\end{longtable}

\section{Operational Policy}

\begin{itemize}
  \item Keep \host{} as an execution backend, not a parallel MATRIX
        development tree.
  \item Keep one adapter per external format.  Gaussian, Molpro and ORCA
        outputs must be normalized before scientific tools consume them.
  \item Use \texttt{gdv32} or \texttt{g16}; do not revive older GDV variants.
  \item Use Molpro 2025.3 through \texttt{module load molpro/2025.3}.
  \item Use ORCA through the configured ORCA module.
  \item Keep credentials outside the repository and outside documentation.
  \item Keep job provenance: input file, MATRIX version, remote command,
        program version, log path and promoted \xyzin{} sections.
\end{itemize}

\end{document}
